\section{Foundations: alphabets, strings, and computation}

\subsection{What is computation?}
A computation is the process by which an abstract machine reads an input
string, updates a finite (or unbounded) internal configuration, and eventually
answers accept or reject. The course studies three models of
increasing memory:
\begin{enumerate}
  \item \textbf{Finite automata} --- finite memory, encoded entirely in the
        current state.
  \item \textbf{Pushdown automata} --- finite control plus a stack.
  \item \textbf{Turing machines} --- unbounded tape memory.
\end{enumerate}
Finite automata correspond to regular languages; pushdown automata to
context-free languages; Turing machines to the recursively enumerable languages.
The first two models occupy these notes; Turing machines come later in the
course.

\subsection{Alphabets, strings, and languages}
\begin{definition}[Alphabet and strings]
An \emph{alphabet} $\Sigma$ is a finite nonempty set of symbols. A
\emph{string} over $\Sigma$ is a finite sequence of symbols from $\Sigma$.
The empty string is $\eps$, and $|\eps|=0$. Write $|w|$ for the length of
$w$. Concatenation of $u$ and $v$ is written $uv$.
\end{definition}

The set of all strings over $\Sigma$ is $\Sigma^*$. Formally
\[
  \Sigma^0=\set{\eps},\qquad
  \Sigma^{n+1}=\set{wa:w\in\Sigma^n,\,a\in\Sigma},\qquad
  \Sigma^*=\bigcup_{n\geq 0}\Sigma^n.
\]
Write $\Sigma^+=\Sigma^*\setminus\set{\eps}$ for the nonempty strings.

\begin{definition}[Language]
A \emph{language} over $\Sigma$ is any subset $L\subseteq\Sigma^*$.
\end{definition}

Two equivalent viewpoints appear throughout the notes.
\begin{itemize}
  \item A machine $M$ that takes strings over $\Sigma$ as input defines
        \[
          \lang{M}=\set{w\in\Sigma^*:M\text{ accepts }w}.
        \]
  \item A function $f:\Sigma^*\to\set{0,1}$ defines
        \[
          L_f=\set{w\in\Sigma^*:f(w)=1}.
        \]
\end{itemize}
Recognition is the special case of computing a Boolean function on strings.

\subsection{First look at finite automata}
A finite automaton is drawn as a \emph{state diagram}:
\begin{itemize}
  \item circles are states;
  \item a double circle is an accept state;
  \item arrows are transitions, labelled by input symbols;
  \item an incoming arrow with no source marks the start state;
  \item the output of a completed run is accept/reject
        (equivalently yes/no).
\end{itemize}

\begin{example}[The machine $M_1$]
Let $\Sigma=\set{0,1}$ and let $M_1$ have states $\set{q_1,q_2,q_3}$, start
state $q_1$, and sole accept state $q_2$.
\begin{center}
\begin{tikzpicture}[automaton,node distance=30mm]
  \node[state,initial] (q1) {$q_1$};
  \node[state,accepting,right=of q1] (q2) {$q_2$};
  \node[state,below=of q2] (q3) {$q_3$};
  \path[->]
    (q1) edge[loop above] node{$0$} ()
         edge node[above]{$1$} (q2)
    (q2) edge[loop above] node{$1$} ()
         edge node[right]{$0$} (q3)
    (q3) edge[bend left=18] node[below left]{$0,1$} (q2);
\end{tikzpicture}
\end{center}
$M_1$ accepts precisely those strings that contain at
least one $1$, after which an even number of $0$s occur. Equivalently: after
the last $1$, the number of trailing $0$s is even (including the case of no
trailing $0$s, i.e.\ strings that end in $1$).

Sample runs:
\[
\begin{array}{c|ccccc}
\text{input} & 1 & 01 & 11 & 010 & 0100\\
\hline
\text{end state} & q_2 & q_2 & q_2 & q_3 & q_2\\
\text{result} & \text{accept} & \text{accept} & \text{accept} & \text{reject} & \text{accept}
\end{array}
\]
The diagram determines a unique language: all $w$ such that
$\transition^*(q_1,w)=q_2$.
\end{example}
